\section{Uniform gains, gradient bounds, and stability}\label{sec:gains}

This section proves the three network estimates that the submatrix theory supplies: a two-sided
per-layer gain bound uniform over all inputs and all gate patterns of a fixed density
(\Cref{thm:uniform-gain}); a lower bound on the gradient that holds simultaneously at every input
(\Cref{prop:gradlb}); and a stability estimate showing the gradient varies slowly, so that a bound at
a net point extends to a neighborhood (\Cref{thm:stability}). Throughout, the network is as in
\Cref{sec:networks} under \Cref{ass:nondeg}, weights are Gaussian with the entry-variance
normalization stated, and the wide regime $dL\log\bar n=o(\min_\ell n_\ell)$ is in force. Proofs of the
lemmas are completed in \Cref{app:netlemmas}.

\subsection{Two-sided per-layer gain}
The forward map of one layer, restricted to a gate of density $\alpha$, gains between
$\sqrt{\kcheck(\alpha)}$ and $\sqrt{\kappa(\alpha)}$ on any direction --- and the extremes are attained
by an adversary choosing which coordinates to keep. This is the network reading of the operator-norm
brackets of \Cref{prop:A2}, made uniform in the direction.

\begin{theorem}[Uniform two-sided gain]\label{thm:uniform-gain}
Let $W$ be a $p\times q$ Gaussian matrix with entry variance $1/q$, and let $V\subseteq\R^q$ be a fixed
subspace with $\dim V=d$. If $d\log p=o(p)$, then w.h.p.\ simultaneously for every $x\in V$ and every
$T\subseteq[p]$ with $|T|=\alpha p$,
\[
  \bigl(\sqrt{\tfrac pq\,\kcheck(\alpha)}-o(1)\bigr)\|x\|\ \le\ \|P_TWx\|\ \le\
  \bigl(\sqrt{\tfrac pq\,\kappa(\alpha)}+o(1)\bigr)\|x\|,
\]
and the extremes are sharp: for each $x$ there are choices of $T$ approaching each end. With the
signal-preserving normalization $p=q$ the factors are $\sqrt{\kcheck(\alpha)}$ and
$\sqrt{\kappa(\alpha)}$.
\end{theorem}

\begin{proof}[Proof sketch; details in \Cref{app:netlemmas}]
For a fixed unit $x\in V$, $Wx$ is a standard Gaussian vector in $\R^p$ scaled by $1/\sqrt q$, so its
coordinates are i.i.d.\ $N(0,1/q)$. The largest $\alpha p$ squared coordinates sum to
$\tfrac pq\kappa(\alpha)(1+o(1))$ and the smallest to $\tfrac pq\kcheck(\alpha)(1+o(1))$ by the
order-statistic law \Cref{lem:order}, with Gaussian concentration giving the deviation probability
$e^{-cp}$ for any fixed threshold. A choice of $T$ selects some $\alpha p$ of the coordinates, so
$\|P_TWx\|^2$ lies between these two extremes for that $x$; taking $T$ to be the top or bottom
coordinates attains them. To make the statement uniform over $x\in V$, take an $\eps$-net of the unit
sphere of $V$, of size $(3/\eps)^d$, apply the bound at each net point with a union bound (the total
failure probability is $(3/\eps)^d e^{-cp}\to0$ since $d\log(1/\eps)=o(p)$), and extend to all $x$ by
the $1$-Lipschitz dependence of $x\mapsto\|P_TWx\|$ on $x$ together with $\|W\|_\op=O(\sqrt{p/q})$; the
supremum over $T$ is handled at the same time because the top/bottom-$\alpha$ selections are themselves
Lipschitz in $x$.
\end{proof}

Applied layer by layer along a network, \Cref{thm:uniform-gain} gives a deterministic two-sided
sandwich for the forward norm and, transposed, for the backward one: uniformly over inputs, every
realized pattern, and the worst adversarial pattern,
\begin{equation}\label{eq:layerbracket}
  \prod_{\ell=1}^L\sqrt{\kcheck(\alpha_\ell)}\ \le\ \frac{\|J(x)v\|}{\|v\|}\ \le\
  \prod_{\ell=1}^L\sqrt{\kappa(\alpha_\ell)}\qquad(v\ \text{worst direction}),
\end{equation}
up to $o(1)$, where $\alpha_\ell$ are the pattern densities. The upper product is the network Lipschitz
constant that the certificate of \Cref{sec:certificates} uses; the lower product is the worst-case
gain, exponentially small in $L$ at the adversarial pattern (each $\kcheck(1/2)=0.071$), which is the
source of the certification gap.

\subsection{Gradient lower bound}
The gradient at a realized input is a specific product at the realized pattern, whose density is
$\Theta(1)$ rather than adversarial, so it does not suffer the $\kcheck$ decay of
\eqref{eq:layerbracket}. With the signal-preserving normalization the backward pass is $\Theta(1)$ at
every input.

\begin{proposition}[Uniform gradient lower bound]\label{prop:gradlb}
Assume the entry variances are chosen signal-preserving, so that $\EE\|h_\ell(x)\|^2=\|h_{\ell-1}(x)\|^2$
(e.g.\ variance $2/n_{\ell-1}$ for ReLU), and $\sigma_\ell$ satisfies the pointwise nondegeneracy
$\sigma_\ell'(r)^2+\sigma_\ell'(-r)^2\ge 2c^2$ for a constant $c>0$ and all $r$. Then there is
$c_g=c_g(c,L)>0$ such that w.h.p.\ over the weights,
\[
  \|\nabla F(x)\|\ \ge\ c_g\,\|w\|\qquad\text{simultaneously for all }x\ \text{on the sphere of radius }R.
\]
For ReLU, $c^2=1/2$ and $c_g$ is bounded below by an absolute constant to the power $L$.
\end{proposition}

\begin{proof}[Proof sketch]
For a fixed $x$, the symmetric-weight decoupling of \cite{bubeckCGT21} replaces each gate $D_\ell(x)$
in distribution by an independent diagonal $\Sigma_\ell$ with entries i.i.d.\ as $\sigma_\ell'(Z)$,
$Z\sim N(0,\cdot)$, so that $\nabla F(x)^\top$ has the law of $w^\top\Sigma_LW_L\cdots\Sigma_1W_1$ with
the $\Sigma$'s independent of the $W$'s. The Frobenius norm of this product obeys the per-layer
recursion $\EE\|\,\cdot\,\|_F^2\ge (2c^2\cdot\tfrac{\text{var factor}}{1})\,\EE\|\,\cdot\,\|_F^2$ of
the previous layer (the pointwise nondegeneracy lower-bounds each layer's contribution and the
signal-preserving variance cancels the mean gate mass), giving
$\EE\|\nabla F(x)\|^2\ge c_g^2\|w\|^2$ with $c_g^2=(2c^2\cdot\tfrac12)^{L}\cdot(1-o(1))=1$ for ReLU up
to lower-order corrections; concentration over the $\Theta(n_L\cdots n_1)$ Gaussian entries gives the
bound whp for that $x$, with failure probability $e^{-c'\min_\ell n_\ell}$. Uniformity over $x$ comes
from an $\eps$-net of the sphere of radius $R$ (size $(CR/\eps)^d$), a union bound (which needs
$d\log(R/\eps)=o(\min_\ell n_\ell)$), and the stability estimate \Cref{thm:stability}, which controls
$\|\nabla F(x')-\nabla F(x)\|$ for $\|x'-x\|\le\eps$. The nondegeneracy hypothesis holds for every
$[0,1]$-slope activation with $\EE[\sigma'(Z)^2]>0$, in particular for ReLU ($c^2=\tfrac12$), leaky
ReLU, and any odd $C^1$ activation with $\sigma'(0)>0$; biases only shift the gate density and are
absorbed into $c$ through \Cref{ass:nondeg}.
\end{proof}

We isolate what is imported and what is new. The fixed-$x$ lower bound is the recursion of
\cite{danielyM26,advExt} together with the decoupling of \cite{bubeckCGT21}; the point here is that it
holds under \Cref{ass:nondeg} with biases and for any $[0,1]$-slope activation, and that its
uniformity over inputs follows from the net plus \Cref{thm:stability}. The stability estimate itself is
where the submatrix operator-norm bracket enters, and it is new in the two-sided form we need.

\subsection{Stability: few gates flip}
Along a short segment the gate pattern changes only at the neurons whose preactivation is near zero,
and there are few of these because Gaussian preactivations anticoncentrate. The Jacobian change is
then an operator-norm perturbation supported on a thin gate, bounded by the max-over-thin-patterns
statistic of \Cref{prop:A2} at the small density of flipped neurons.

\begin{theorem}[Gradient stability]\label{thm:stability}
Under the hypotheses of \Cref{prop:gradlb}, there is a constant $C=C(L)$ such that w.h.p.\ over the
weights, for all inputs $x$ on the sphere and all $x'$ with $\|x'-x\|\le r$,
\[
  \|\nabla F(x')-\nabla F(x)\|\ \le\ C\,\|w\|\,\Bigl(\tfrac{r\sqrt{\bar n}}{\|x\|}\log\tfrac{\|x\|}{r\sqrt{\bar n}}\Bigr)^{1/2}.
\]
In particular, for $r\le \|x\|/\sqrt{\bar n}\cdot\eta$ the right side is $C\|w\|(\eta\log(1/\eta))^{1/2}$,
which is $o(\|w\|)$ as $\eta\to0$; the stability radius needed for the attack, $r\asymp\|x\|/\sqrt{n_L}$
up to logarithms, gives a change small compared with the lower bound $c_g\|w\|$ of \Cref{prop:gradlb}.
\end{theorem}

\begin{proof}[Proof sketch; details in \Cref{app:netlemmas}]
Write the Jacobian difference by the telescoping identity
\[
  \nabla F(x')^\top-\nabla F(x)^\top=\sum_{\ell=1}^L
  w^\top D_L(x')W_L\cdots\bigl(D_\ell(x')-D_\ell(x)\bigr)W_\ell\cdots D_1(x)W_1,
\]
so it suffices to bound each $\|(D_\ell(x')-D_\ell(x))W_\ell(\cdots)\|$. The gate difference
$D_\ell(x')-D_\ell(x)$ is supported on the neurons whose slope changed between $x$ and $x'$, i.e.\ those
with a preactivation sign change; a neuron flips only if its preactivation at $x$ is within
$O(\|z_\ell(x')-z_\ell(x)\|_\infty)$ of zero. Since the preactivations are Gaussian of scale
$\varsigma_\ell\asymp\|h_{\ell-1}\|/\sqrt{n_{\ell-1}}\asymp\|x\|/\sqrt{n_{\ell-1}}$ and the layer map is
$O(1)$-Lipschitz, the number of flipped neurons at layer $\ell$ is, by the anticoncentration bound
$\PP(|N(0,\varsigma^2)|\le\tau)\le\tau/\varsigma$, at most $n_\ell\cdot O(r\sqrt{n_{\ell-1}}/\|x\|)=:%
\delta_\ell n_\ell$ w.h.p., with $\delta_\ell=O(r\sqrt{\bar n}/\|x\|)$. The factor
$(D_\ell(x')-D_\ell(x))W_\ell$ is then a Gaussian block on a pattern of density $\delta_\ell$, whose
operator norm is at most $2\sqrt{\kappa(\delta_\ell)}+o(1)$ by \Cref{prop:A2}, and
$\kappa(\delta)\le C\delta\log(1/\delta)$ for small $\delta$ (the top-$\delta$ mass of a Gaussian).
The remaining factors are $O(1)$ in operator norm by \Cref{thm:uniform-gain}, and $\|w^\top D_L\cdots\|$
contributes $\|w\|$; multiplying and summing over $\ell$ gives the stated bound, with the logarithm
from $\kappa(\delta_\ell)\le C\delta_\ell\log(1/\delta_\ell)$. Uniformity over $x,x'$ is again by a net
plus the $O(1)$-Lipschitz dependence of every factor.
\end{proof}

The three estimates combine into the input of the constructive attack: a small output (next section),
a gradient bounded below by $c_g\|w\|$ (\Cref{prop:gradlb}), and a gradient stable to within
$o(\|w\|)$ over the attack segment (\Cref{thm:stability}). The compatibility of these three is what
\Cref{sec:fragility} verifies.
